\documentclass[a4paper]{article}

%====================== PACKAGES ======================

\usepackage[french]{babel}
\usepackage[utf8x]{inputenc}
%pour gérer les positionnement d'images
\usepackage{float}
\usepackage{amsmath}
\usepackage{graphicx}
\usepackage[colorinlistoftodos]{todonotes}
\usepackage{url}
%pour les informations sur un document compilé en PDF et les liens externes / internes
\usepackage{hyperref}
%pour la mise en page des tableaux
\usepackage{array}
\usepackage{tabularx}
%pour utiliser \floatbarrier
%\usepackage{placeins}
%\usepackage{floatrow}
%espacement entre les lignes
\usepackage{setspace}
%modifier la mise en page de l'abstract
\usepackage{abstract}
%police et mise en page (marges) du document
\usepackage[T1]{fontenc}
\usepackage[top=2cm, bottom=2cm, left=2cm, right=2cm]{geometry}
%Pour les galerie d'images
\usepackage{subfig}

%====================== AFFICHAGE DE LISTINGS ======================

\usepackage{listings}

\lstset{
	language=Python,
	basicstyle=\ttm,
	otherkeywords={self},             % Add keywords here
	keywordstyle=\ttb\color{deepblue},
	emph={MyClass,__init__},          % Custom highlighting
	emphstyle=\ttb\color{deepred},    % Custom highlighting style
	stringstyle=\color{deepgreen},
	frame=tb,                         % Any extra options here
	showstringspaces=false            %
}

\definecolor{javared}{rgb}{0.6,0,0} % for strings
\definecolor{javagreen}{rgb}{0.25,0.5,0.35} % comments
\definecolor{javapurple}{rgb}{0.5,0,0.35} % keywords
\definecolor{javadocblue}{rgb}{0.25,0.35,0.75} % javadoc

\lstset{language=Java,
	basicstyle=\ttfamily\footnotesize,
	frame=single,
	keywordstyle=\color{javapurple}\bfseries,
	stringstyle=\color{javared},
	commentstyle=\color{javagreen},
	morecomment=[s][\color{javadocblue}]{/**}{*/},
	numbers=left,
	numberstyle=\tiny\color{black},
	stepnumber=2,
	numbersep=10pt,
	breaklines=true,
	tabsize=2,
	showspaces=false,
	showstringspaces=false}

%====================== INFORMATION ET REGLES ======================

%rajouter les numérotation pour les \paragraphe et \subparagraphe
\setcounter{secnumdepth}{4}
\setcounter{tocdepth}{4}

\hypersetup{							% Information sur le document
pdfauthor = {Premier Auteur,
			Deuxième Auteur,
			Troisième Auteur,
    		Quatrième Auteur},			% Auteurs
pdftitle = {Nom du Projet -
			Sujet du Projet},			% Titre du document
pdfsubject = {Mémoire de Projet},		% Sujet
pdfkeywords = {Tag1, Tag2, Tag3, ...},	% Mots-clefs
pdfstartview={FitH}}					% ajuste la page à la largueur de l'écran
%pdfcreator = {MikTeX},% Logiciel qui a crée le document
%pdfproducer = {}} % Société avec produit le logiciel

%======================== DEBUT DU DOCUMENT ========================

\begin{document}

%régler l'espacement entre les lignes
\newcommand{\HRule}{\rule{\linewidth}{0.5mm}}

%page de garde
\begin{titlepage}
\begin{center}

% le tilde a la ligne suivante est un espace insécable; il a pour effet de commencer
% un paragraphe. Si il n'était pas là, la commande d'espacement bigskip n'aurait pas 
% d'effet. On peut mettre plus de commandes d'espacement pour descendre le logo plus bas,
% si souhaité.
~\bigskip

\includegraphics[width=0.35\textwidth]{./logo}\\[1cm]

\textsc{\LARGE Université ou Entreprise}\\[1.5cm]

% Title
\HRule \\[0.4cm]

{\huge \bfseries Nom du Projet\\
Sujet du Projet \\[0.4cm] }

\HRule \\[1.5cm]

% Author and supervisor
\begin{minipage}{0.4\textwidth}
\begin{flushleft} \large
\emph{Auteurs:}\\
Premier \textsc{Auteur}\\
Deuxième \textsc{Auteur}\\
Troisième \textsc{Auteur}\\
Quatrième \textsc{Auteur}
\end{flushleft}
\end{minipage}
\begin{minipage}{0.4\textwidth}
\begin{flushright} \large
\emph{Client:} \\
Prénom \textsc{Nom}\\
\emph{Référent:} \\
Prénom \textsc{Nom}
\end{flushright}
\end{minipage}

~\bigskip
\bigskip
\bigskip
\bigskip

\begin{abstract}

\begin{spacing}{1.3}

Lorem ipsum dolor sit amet, consectetur adipiscing elit. Sed non risus. Suspendisse lectus tortor, dignissim sit amet, adipiscing nec, ultricies sed, dolor. Cras elementum ultrices diam. Maecenas ligula massa, varius a, semper congue, euismod non, mi. Proin porttitor, orci nec nonummy molestie, enim est eleifend mi, non fermentum diam nisl sit amet erat. Duis semper. Duis arcu massa, scelerisque vitae, consequat in, pretium a, enim. Pellentesque congue. Ut in risus volutpat libero pharetra tempor. Cras vestibulum bibendum augue. Praesent egestas leo in pede. Praesent blandit odio eu enim. Pellentesque sed dui ut augue blandit sodales. Vestibulum ante ipsum primis in faucibus orci luctus et ultrices posuere cubilia Curae; Aliquam nibh. Mauris ac mauris sed pede pellentesque fermentum. Maecenas adipiscing ante non diam sodales hendrerit. Ut velit mauris, egestas sed, gravida nec, ornare ut, mi. Aenean ut orci vel massa suscipit pulvinar. Nulla sollicitudin. Fusce varius, ligula non tempus aliquam, nunc turpis ullamcorper nibh, in tempus sapien eros vitae ligula. Pellentesque rhoncus nunc et augue. Integer id felis. Curabitur aliquet pellentesque diam. 

\end{spacing}
\end{abstract}

% Afficher la date de compilation en bas de la page de titre
% On peut remplacer \today par une date fixe, ou commenter la ligne
% si on ne veut rien afficher du tout.
\vfill
{\large \today}

\end{center}
\end{titlepage}


%%% si on souhaite une table des matières, on peut dé-commenter les lignes
%%% suivantes (de { a })
% \tableofcontents
% \thispagestyle{empty}
% \setcounter{page}{0}
% %ne pas numéroter le sommaire
%
% \newpage

%espacement entre les lignes d'un tableau
\renewcommand{\arraystretch}{1.5}

%recommencer la numérotation des pages à "1"
\setcounter{page}{1}

% le contenu du rapport va dans cette partie
\section{Premiere section}


% On préfère commencer la bibliographie sur une nouvelle page, 
% mais on peut commenter la commande suivante pour qu'elle suive
% immédiatement la fin du dernier chapitre
\newpage

%choix du style de la bibliographie
\bibliographystyle{plain}
%inclusion de la biblio
\bibliography{bibliographie.bib}
%voir wiki pour plus d'information sur la syntaxe des entrées d'une bibliographie

\end{document}
